La ciencia de datos puede ampliar el acceso a información productiva allí donde los registros de cosecha son incompletos o costosos de obtener. Las estimaciones parcelarias de rendimiento permiten complementar evaluaciones de capacidad productiva, planificación de flujos de efectivo, seguimiento territorial del riesgo y selección de necesidades de asistencia técnica. La observación satelital aporta indicadores comparables de vegetación, sequía y productividad hídrica cuya obtención no exige visitar simultáneamente todas las parcelas \parencite{Bokusheva2016,Dhonthi2024,FAOWaPOR,Serban2025}.

En el análisis financiero, sin embargo, rendimiento no equivale a rentabilidad, ingreso neto ni solvencia. Para valorar la capacidad de pago deben considerarse superficie realmente cosechada, precio de venta, costos de producción, exposición climática, otros ingresos y compromisos crediticios. Tampoco es lícito extrapolar sin reservas los indicadores de 138 parcelas a todos los productores de Hidalgo, Puebla y Tlaxcala; sus poblaciones, sistemas de producción y mecanismos de selección pueden ser diferentes de los contenidos en las estadísticas oficiales \parencite{FIRA2026AgroCebada,SIAPDatos}.

Una aplicación responsable requeriría conservar procedencia y fecha de cada predictor, revisar calidad de las observaciones, comprobar estabilidad territorial, cuantificar incertidumbre con información independiente y habilitar corrección humana ante situaciones atípicas. Las coordenadas y los rendimientos individuales deben protegerse por su sensibilidad comercial. Las estimaciones no deberían utilizarse como regla automática y exclusiva para aprobar o denegar un crédito, pues podrían desfavorecer sistemáticamente a parcelas o grupos insuficientemente representados.

En consecuencia, la principal contribución a la inclusión financiera no es sustituir la evaluación agronómica por una cifra algorítmica, sino ofrecer un insumo adicional, trazable y sometido a auditoría. La reducción de asimetrías informativas solamente genera valor sostenible cuando se acompaña de mecanismos que reconozcan explícitamente el error de predicción, el contexto productivo y las limitaciones de transferencia a otros ciclos.
